% ej_coord_R2_02.tex
%
% Copyright (C) 2022--2026 José A. Navarro Ramón <janr.devel@gmail.com>
% Licencia del código GPLv2
% Licencia Creative Commons Recognition Non-Commercial Share-alike.
% (CC-BY-NC-SA)

\section{Ejercicio 2}
Supongamos un cierto vector en coordenadas rectangulares,
$\vvv{v} = 2\,\uvec{\i} + \sqrt{3}\,\uvec{\j}$.\\
a) Calcula el cuadrado del módulo de $\vvv{v}$.\\
b) Queremos utilizar un sistema de coordenadas oblicuo, cuyos ejes forman
$\qty{60}{\degree}$ entre sí y haciendo coincidir los ejes $x$ y $x'$.
Calcula los vectores de la base $B'=\set{\vvv{e}_1, \vvv{e}_2}$ en el sistema oblicuo
en función de la base del sistema cartesiano ortogonal $B=\set{\uvec{\i},\,\uvec{\j}}$.
Halla sus módulos y comprueba que son vectores unitarios.\\
c) Halla la matriz de cambio de base del sistema cartesiano ortogonal
$B=\set{\uvec{\i},\,\uvec{\j}}$ a la del sistema de coordenadas oblicuas, $B'$.
Calcula también su inversa.\\
d) Obtén las componentes contravariantes del vector $\vvv{v}$ en el sistema de
coordenadas oblicuas, mediante una transformación de componentes.\\
e) Halla ahora las componentes, mediante transformación de los vectores de la base $B$.\\
f) Calcula el producto escalar de los vectores de la base covariante
$B'=\set{\vvv{e}_1, \vvv{e}_2}$ en función de la base en coordenadas ortogonales.\\
g) Obtén el tensor métrico en el sistema oblicuo y su inversa.\\
h) Calcula los vectores de la base contravariante $\set{\vvv{e}^1,\vvv{e}^2}$ en función
de la base rectangular $\set{\uvec{\i},\,\uvec{\j}}$ y halla sus módulos, comprobando
que no son vectores unitarios\\
i) Halla el producto escalar entre los elementos de la base covariante y contravariante.\\
j) Halla las componentes covariantes del vector.\\
k) Halla el cuadrado del módulo del vector en el sistema oblicuo.\\
l) Relaciona el apartado anterior, con el hecho de que los vectores covariantes son
elementos $\omega$ del espacio dual de los vectores $\vvv{v}$ del espacio tangente de $\symbb{R}^2$, y que estas $\omega$ son aplicaciones lineales que al actuar sobre los
vectores del espacio tangente, producen un escalar. Halla el escalar para nuestro vector
$\vvv{v}$.\\

\noindent
\textbf{Solución}

\noindent
a) Cuadrado del módulo de $\vvv{v} = 2\,\uvec{\i} + \sqrt{3}\,\uvec{\j}$
\[
  |\vvv{v}|^2 = v^2 = v_x^2 + v_y^2 = 2^2 + \left(\sqrt{3}\right)^2 = 4 + 3 = 7
\]

\noindent 
b) Matriz de cambio de base de $B=\set{\uvec{\i},\,\uvec{\j}}$ a la de coordenadas
oblicuas que forman $\qty{60}{\degree}$, haciendo coincidir los ejes $x$ y $x'$.
Y también su inversa.

En la figura \ref{fig:R2-coord_oblicuas_e1e2} se puede observar la relación entre estos
vectores, y además podemos comprobar que son unitarios
\begin{align*}
  \vvv{e}_1
  &=
    \uvec{\i}
    \longrightarrow
    |\vvv{e}_1| = |\uvec{\i}|
    = 1\\
  \vvv{e}_2
  &=
    \cos(\pi/3)\,\uvec{\i} + \sin(\pi/3)\,\uvec{\j}
    = \dfrac{1}{2}\,\uvec{\i} + \dfrac{\sqrt{3}}{2}\,\uvec{\j}
    \longrightarrow
    |\vvv{e}_2| = \left|\dfrac{1}{2}\,\uvec{\i}+\dfrac{\sqrt{3}}{2}\,\uvec{\j}\right|
    = \dfrac{1}{4} + \dfrac{3}{4}
    = 1
\end{align*}

\noindent
c) Matriz de cambio de base $B$ a $B'$ y su inversa.

La transformación de vectores de la base es $B'=B \mmm{M}$.
Del apartado anterior se puede comprobar que
\[
  \begin{pmatrix}
    \vvv{e}_1 & \vvv{e}_2
  \end{pmatrix}
  = \begin{pmatrix}
    \uvec{\i} & \uvec{\j}
  \end{pmatrix}
  \,\begin{pmatrix}
    1 & 1/2\\
    0 & \sqrt{3}/2
  \end{pmatrix}
\]

Luego, la matriz de cambio de base de $B$ a $B'$ es
\[
  \mmm{M}
  = \begin{pmatrix}
    1 & 1/2\\
    0 & \sqrt{3}/2
  \end{pmatrix}
\]

Su inversa es
\[
  \mmm{M}^{-1}
  = \dfrac{1}{\det(\mmm{M})}\,
  \text{adj}(\mmm{M}^\transp)
  = \dfrac{1}{\sqrt{3}/2}
  \begin{pmatrix}
    \sqrt{3}/2 & -1/2\\
    0 & 1
  \end{pmatrix}
  = \begin{pmatrix}
    1 & -\sqrt{3}/3\\
    0 & 2\sqrt{3}/3
  \end{pmatrix}
\]

\noindent
d) Componentes contravariantes de $\vvv{v}$.

Estas componentes se calculan con la inversa de la matriz de cambio de base
\[
  \begin{pmatrix}
    x^1\\
    x^2
  \end{pmatrix}
  = \begin{pmatrix}
    1 & -\sqrt{3}/3\\
    0 & 2\sqrt{3}/3
  \end{pmatrix}
  \,\begin{pmatrix}
    x\\
    y
  \end{pmatrix}
\]
\[
  \begin{pmatrix}
    x^1\\
    x^2
  \end{pmatrix}
  = \begin{pmatrix}
    1 & -\sqrt{3}/3\\
    0 & 2\sqrt{3}/3
  \end{pmatrix}
  \,\begin{pmatrix}
    2\\[0.4ex]
    \sqrt{3}
  \end{pmatrix}
  = \begin{pmatrix}
    1\\
    2
  \end{pmatrix}
\]

Las coordenadas contravariantes en el nuevo sistema son
\begin{align*}
  x^1 = 1\\
  x^2 = 2
\end{align*}

\noindent
e) Componentes contravariantes mediante transformación de vectores de la base $B$.

Necesitamos obtener los vectores de la base antigua $B=\set{\uvec{\i},\uvec{\j}}$ en
función de la nueva $B'=\set{\vvv{e}_1,\vvv{e}_2}$. Para ello utilizaremos la inversa de
la matriz de cambio de base
\[
  \begin{pmatrix}
    \uvec{\i} & \uvec{\j}
  \end{pmatrix}
  =
  \begin{pmatrix}
    \vvv{e}_1 & \vvv{e}_2
  \end{pmatrix}
  \begin{pmatrix}
    1 & -\dfrac{\sqrt{3}}{3}\\[1.4ex]
    0 & \dfrac{2\sqrt{3}}{3}
  \end{pmatrix}
  =
  \begin{pmatrix}
    \vvv{e}_1 & -\dfrac{\sqrt{3}}{3}\vvv{e}_1+\dfrac{2\sqrt{3}}{3}\vvv{e}_2
  \end{pmatrix}
\]

La transformación de vectores resulta
\begin{align*}
  \uvec{\i} &= \vvv{e}_1\\
  \uvec{\j} &= -\dfrac{\sqrt{3}}{3}\vvv{e}_1 + \dfrac{2\sqrt{3}}{3}\vvv{e}_2
\end{align*}

Por último, se sustituyen estos vectores de $B$ en la expresión de $\vvv{v}$
\[
  \vvv{v}
  =
  2\,\uvec{\i} + \sqrt{3}\,\uvec{\j}
  = 2 \vvv{e}_1
  +\sqrt{3}\left(-\dfrac{\sqrt{3}}{3}\vvv{e}_1+\dfrac{2\sqrt{3}}{3}\vvv{e}_2\right)
  =
  2\vvv{e}_1 - \vvv{e}_1 + 2\vvv{e}_2
  = \vvv{e}_1 + 2\vvv{e}_2 
\]

El vector expresado con componentes contravariantes es
\[
  \vvv{v} = x^i\vvv{e}_i = x^1\vvv{e}_1 + x^2\vvv{e}_2
\]

Comparando con nuestro resultado
\[
  \vvv{v} = x^1\vvv{e}_1 + x^2\vvv{e}_2 = 1\vvv{e}_1 + 2\vvv{e}_2
\]

Estas componentes son, obviamente, las mismas que en el apartado anterior
\begin{align*}
  x^1 &= 1\\
  x^2 &= 2
\end{align*}

\noindent
f) Producto escalar de los vectores de la base covariante $B=\set{\vvv{e}_1,\vvv{e}_2}$.
\begin{align*}
  \vvv{e}_1\cdot\vvv{e}_1
  &=
    \uvec{\i}\cdot\uvec{\i}
    = 1\\
  \vvv{e}_1\cdot\vvv{e}_2
  &=
    \vvv{e}_2\cdot\vvv{e}_1
    = \uvec{\i}\cdot\left(\dfrac{1}{2}\,\uvec{\i}+\dfrac{\sqrt{3}}{2}\,\uvec{\j}\right)
    = \dfrac{1}{2}\\
  \vvv{e}_2\cdot\vvv{e}_2
  &=
    \left(\dfrac{1}{2}\,\uvec{\i}+\dfrac{\sqrt{3}}{2}\,\uvec{\j}\right)
    \cdot
    \left(\dfrac{1}{2}\,\uvec{\i}+\dfrac{\sqrt{3}}{2}\,\uvec{\j}\right)
    = \left(\dfrac{1}{2}\right)^2 + \left(\dfrac{\sqrt{3}}{2}\right)^2
    = \dfrac{1}{4} + \dfrac{3}{4}
    = 1
\end{align*}

\noindent
g) Tensor métrico y su inverso.

El tensor métrico sería
\[
  \mmm{g}_{ij}
  =
  \begin{pmatrix}
    g_{11} & g_{12}\\
    g_{21} & g_{22}
  \end{pmatrix}
  =
  \begin{pmatrix}
    \xhat{e}_1\cdot\xhat{e}_1 & \xhat{e}_1\cdot\xhat{e}_2\\
    \xhat{e}_2\cdot\xhat{e}_1 & \xhat{e}_2\cdot\xhat{e}_2
  \end{pmatrix}
  =
  \begin{pmatrix}
    1 & 1/2\\
    1/2 & 1
  \end{pmatrix}
\]

La inversa del tensor métrico es
\[
  \mmm{g}^{ij}
  = \dfrac{1}{\det(\mmm{g}_{ij})}\, \text{adj}(\mmm{g}_{ij}^\transp)
  = \dfrac{1}{3/4}
  \begin{pmatrix}
    1 & -1/2\\
    -1/2 & 1
  \end{pmatrix}
  = \begin{pmatrix}
    4/3 & -2/3\\
    -2/3 & 4/3
  \end{pmatrix}
\]

\noindent
h) Base contravariante en función de $\uvec{\i}$ y $\uvec{\j}$.

En el apartado d) se calcularon los vectores de la base covariante, resultando
\begin{align*}
  \vvv{e}_1 &= \uvec{\i}\\
  \vvv{e}_2 &= \dfrac{1}{2}\,\uvec{\i} + \dfrac{\sqrt{3}}{2}\,\uvec{\j}
\end{align*}

Subiendo índices obtendríamos los vectores de la base contravariante
\[
  \vvv{e}^i = g^{ij}\vvv{e}_j
\]
\begin{align*}
  \vvv{e}^1
  &=
    g^{11}\vvv{e}_1 + g^{12}\vvv{e}_2
    = \dfrac{4}{3}\vvv{e}_1 - \dfrac{2}{3}\vvv{e}_2
    = \dfrac{4}{3}\,\uvec{\i}
    - \dfrac{2}{3}\,\left(\dfrac{1}{2}\,\uvec{\i}+\dfrac{\sqrt{3}}{2}\,\uvec{\j}\right)
    = \uvec{\i} - \dfrac{\sqrt{3}}{3}\,\uvec{\j}\\
  \vvv{e}^2
  &=
    g^{21}\vvv{e}_1 + g^{22}\vvv{e}_2
    = -\dfrac{2}{3}\vvv{e}_1 + \dfrac{4}{3}\vvv{e}_2
    = -\dfrac{2}{3}\,\uvec{\i}
    + \dfrac{4}{3}\,\left(\dfrac{1}{2}\,\uvec{\i}+\dfrac{\sqrt{3}}{2}\,\uvec{\j}\right)
    = \dfrac{2\sqrt{3}}{3}\,\uvec{\j}  
\end{align*}

Los vectores de la base contravariante son
\begin{align*}
  \vvv{e}^1 &= \uvec{\i} - \dfrac{\sqrt{3}}{3}\,\uvec{\j}\\
  \vvv{e}^2 &= \dfrac{2\sqrt{3}}{3}\,\uvec{\j}
\end{align*}

Calculamos sus módulos y comprobamos que no son vectores unitarios
\begin{align*}
  \left|\vvv{e}^1\right|
  &=
    1^2 + \left(\dfrac{\sqrt{3}}{3}\right)^2
    = 1 + \dfrac{1}{3}
    = \dfrac{4}{3}\\
  \left|\vvv{e}^2\right|
  &=
    \left(\dfrac{2\sqrt{3}}{3}\right)^2
    = \dfrac{4}{3}
\end{align*}

\noindent
i) Producto escalar entre elementos de las bases covariante y contravariante
\begin{align*}
  \vvv{e}^1\cdot\vvv{e}_1
  &=
    \left(\uvec{\i} - \dfrac{\sqrt{3}}{3}\,\uvec{\j}\right)
    \cdot
    \uvec{\i}
    = 1\\
  \vvv{e}^2\cdot\vvv{e}_2
  &=
    \left(\dfrac{2\sqrt{3}}{3}\,\uvec{\j}\right)
    \cdot
    \left(\dfrac{1}{2}\,\uvec{\i} + \dfrac{\sqrt{3}}{2}\,\uvec{\j}\right)
    = 1
\end{align*}
Y los productos cruzados
\begin{align*}
  \vvv{e}^1\cdot\vvv{e}_2
  &=
    \left(\uvec{\i} - \dfrac{\sqrt{3}}{3}\,\uvec{\j}\right)
    \cdot
    \left(\dfrac{1}{2}\,\uvec{\i} + \dfrac{\sqrt{3}}{2}\,\uvec{\j}\right)
    = \dfrac{1}{2} - \dfrac{\sqrt{3}}{3}\cdot\dfrac{\sqrt{3}}{2}
    = \dfrac{1}{2} - \dfrac{1}{2}
    = 0\\
  \vvv{e}^2\cdot\vvv{e}_1
  &=
    \left(
    \dfrac{2\sqrt{3}}{3}\,\uvec{\j}\right) \cdot \uvec{\i}
    = 0
\end{align*}

\noindent
j) Componentes covariantes de $\vvv{v}$.

Las componentes covariantes del vector se pueden calcular bajando el índice de las
contravariantes
\[
  x_i = g_{ij} x^j
\]
\begin{align*}
  x_1
  &=
    g_{11} x^1 + g_{12} x^2
    = 1\cdot 1 + \dfrac{1}{2} \cdot 2
    =  2\\
  x_2
  &=
    g_{21} x^1 + g_{22} x^2
    = \dfrac{1}{2} \cdot 1 + 1\cdot 2
    = \dfrac{5}{2}
\end{align*}

\noindent
k) Cuadrado de $\vvv{v}$ en el sistema oblicuo.

Se comprueba que el vector no ha cambiado. Solo lo han hecho sus componentes en el nuevo
sistema de coordenadas
\[
  |\vvv{v}|^2
  = v^2
  = x_ix^i
  = x_1x^1 + x_2x^2
  = 2\cdot 1 + \dfrac{5}{2}\cdot 2
  = 2 + 5
  = 7
\]

\noindent
l) Explica el apartado anterior, con el hecho de que los vectores covariantes son
elementos $\omega$ del espacio dual de los vectores $\vvv{v}$ del espacio tangente de
$\symbb{R}^2$ y que estas $\omega$ son aplicaciones lineales que al actuar sobre un
vector, producen un escalar.

El vector $\vvv{v}$ es un elemento del espacio tangente de $\symbb{R}^2$
\[
  \vvv{v} \in T_p\symbb{R}^2
\]

Podemos establecer una aplicación lineal $\omega$ que, al actuar sobre un vector, produce
un escalar
\[
  \omega: T_p\symbb{R}^2 \rightarrow \symbb{R}
\]

Diremos que $\omega$ es un elemento del \emph{espacio dual} del vector
\[
  \omega \in \left(T_p\symbb{R}^2\right)^*
\]

Los vectores covariantes son estas aplicaciones lineales $\omega = x_ie^i$ y, como hemos
dicho, son elementos del espacio dual, que no es el espacio tangente donde \emph{viven}
los vectores. La métrica permite relacionar los dos espacios (subiendo o bajando índices).

Ahora nos podemos preguntar que si
$\omega = x_1\vvv{e}^1 + x_2\vvv{e}^2 = 2\vvv{e}^1 + \frac{5}{2}\vvv{e}^2$, es una de
tales aplicaciones lineales, debería \emph{comer} un vector como $\vvv{v}$ y producir un
escalar (en este caso un número real).

Veamos en notación de Einstein, que el número real es el cuadrado del módulo del vector
\[
  \omega(\vvv{v}) = x_ie^i\cdot x^ie_i = x_ix^i = v^2
\]

En nuestro ejemplo, tendríamos
\[
  \omega(\vvv{v})
  = \left(2\vvv{e}^1+\dfrac{5}{2}\vvv{e}^2\right)\cdot\left(1\vvv{e}_1+2\vvv{e}_2\right)
  = 2\cdot 1 + \dfrac{5}{2}\cdot 2
  = 2 + 5
  = 7
\]
donde se han tenido en cuenta los productos escalares del apartado anterior.

Recalcamos el hecho de que la aplicación lineal $\omega\in (T_p\symbb{R}^2)^*$
(elemento del espacio dual de $\vvv{v} = 2\,\uvec{\i}+\sqrt{3}\,\uvec{\j}$), que es
$\omega=2\vvv{e}^1+\frac{5}{2}\vvv{e}^2$, cuando se aplica al vector $\vvv{v}$ nos da un
escalar, que resulta ser el cuadrado del módulo del vector



%%% Local Variables:
%%% mode: latex
%%% TeX-engine: luatex
%%% TeX-master: "../retazosmatematicas.tex"
%%% End:
